\section{Short transforms and band reduction}
\label{intro:section}

Let \(v:\R^2\to S^1\) be a Lipschitz unit vector field. For
\(0<\epsilon<a\), define the short directional Hilbert transform by
\[
 H_{v,a}^{\epsilon}f(x)
 =\int_{\epsilon<|t|<a}f(x-tv(x))\,\frac{\mathrm dt}{t}.
\]
The direction is held fixed at the output point \(x\): the integral
samples the straight line through \(x\) in direction \(v(x)\).
For \(f\in\mathcal S(\R^2)\), the symmetric principal value
\(H_{v,a}f=\lim_{\epsilon\downarrow0}H_{v,a}^{\epsilon}f\)
exists at every point, since
\(f(x-tv(x))-f(x+tv(x))=O_x(|t|)\) near zero.

\begin{theorem}\label{thm:main}
There are absolute constants \(a_*\in(0,1/2)\) and \(C_*<\infty\)
such that every Lipschitz field \(v:\R^2\to S^1\) with
\(\Lip(v)\le1\) satisfies
\[
 \sup_{0<\epsilon<a_*}
 \norm{H_{v,a_*}^{\epsilon}f}_{L^2(\R^2)}
 \le C_*\norm{f}_{L^2(\R^2)}
 \qquad(f\in\mathcal S(\R^2)).
\]
Its symmetric principal value therefore satisfies
\[
 \norm{H_{v,a_*}f}_2\le C_*\norm f_2,
 \qquad
 \big|\{x:|H_{v,a_*}f(x)|>\lambda\}\big|
 \le C_*^2\lambda^{-2}\norm f_2^2
 \quad(\lambda>0).
\]
\end{theorem}

Both coordinates may enter the field arbitrarily subject to the
Lipschitz bound. The supremum in the theorem is outside the norm,
so its first assertion is a uniform operator bound for each inner
truncation. The principal-value assertion is stated on Schwartz
functions and gives an \(L^2\)-bounded extension by density.

The outer length has a natural scaling. If \(L_v=\Lip(v)>0\), put
\(\widetilde v(X)=v(X/L_v)\) and
\(\widetilde f(X)=f(X/L_v)\). Then \(\Lip(\widetilde v)=1\), and
changing variables \(s=L_vt\) gives
\[
 H_{v,a_*/L_v}^{\epsilon}f(x)
 =H_{\widetilde v,a_*}^{L_v\epsilon}\widetilde f(L_vx),
 \qquad \norm{\widetilde f}_2=L_v\norm f_2.
\]
The theorem thus has the same constant at outer length \(a_*/L_v\).
If \(L_v=0\), the field is constant. Rotation and Fubini reduce its
transform, for any finite outer length, to the one-dimensional
Hilbert transform. These conventions distinguish the short-scale
assertion from an estimate at arbitrary outer lengths for varying
fields.

\subsection{The problem and its predecessors}

Stein's question places Lipschitz regularity at the center of the
planar directional singular-integral problem. In the formulation
\cite[Conjecture~1.5]{LLStein}, it asks for a uniform weak-\((2,2)\)
bound with the outer length controlled by the reciprocal Lipschitz
seminorm. \Cref{thm:main} gives an affirmative answer with a strong
\(L^2\) bound. The scale restriction limits the turning of the field,
but it does not make the operator a Fourier multiplier: the line
used by the integral still varies with the output point.

Lacey and Li trace Stein's question to Zygmund's conjecture on
almost-everywhere differentiation by averages over shrinking straight
segments in direction \(v(x)\). They also recall a Besicovitch
obstruction showing that H\"older regularity of any exponent below one
does not by itself give the corresponding uniform weak-\((2,2)\)
estimate for the short positive averaging maximal operator
\cite[Preface, Conjecture~1.2 and Chapter~2]{LLStein}.
Analytic fields admit additional structure:
Bourgain proved an \(L^2\) theorem for their averaging maximal
operator \cite{Bourgain1989}, and Stein and Street proved local
\(L^p\), \(1<p<\infty\), estimates for single-parameter singular
Radon transforms associated with real-analytic maps
\cite{SteinStreet2012}. Their
singular-integral theorem includes analytic frozen-direction Hilbert
transforms, with localization depending on the analytic map. The
uniform Lipschitz estimate considered here asks for constants
controlled by the Lipschitz seminorm alone.

Several distinct advances explain the remaining difficulty. Lacey
and Li proved an annular weak-\((2,2)\) estimate, together with strong
\(L^p\) estimates for \(p>2\), for directional Hilbert transforms
with measurable directions \cite[Theorem~1.1]{LLAnnular}. An annular
restriction isolates a frequency scale; it does not give a uniform
strong \(L^2\) bound for the complete transform. For Lipschitz unit
fields, Guo obtained annular weak-\((2,2)\) and strong \(L^p\),
\(p>2\), bounds, uniformly in the annulus scale, for the
\(r\)-variation, \(r>2\), of
smooth dyadic Hilbert-transform partial sums over the short scale range
determined by the reciprocal Lipschitz length; the varying dyadic
cutoff is distinct from the hard inner time cutoff in \Cref{thm:main}
\cite[Theorem~1.2 and (1.6)--(1.8) of the author version]{GuoVariation2017}.
Lacey and Li's Lipschitz Kakeya theorem controls averages over short
rectangles on which the field has a prescribed popular direction
\cite[Theorem~1.4]{LLKakeya}.
For popularity \(\delta\), its weak-\(L^2\) norm costs
\(C\delta^{-1/2}\). We use precisely this estimate.

The Lacey--Li memoir develops the relation between those maximal
operators and the Hilbert-transform problem \cite{LLStein}. Its
Theorem~1.18 assumes a stronger maximal estimate below \(L^2\) for
the strong annular conclusion; its full-transform conclusion also
requires \(C^{1+\eta}\) regularity. These are conditional results with
hypotheses different from those of \Cref{thm:main}. For bounded
measurable slopes depending on one coordinate, Bateman proved
single-annulus \(L^p\) estimates for \(1<p<\infty\)
\cite{Bateman2013}. Bateman and Thiele proved strong \(L^p\) bounds
for the full-line principal-value transform when \(3/2<p<\infty\),
for every nonvanishing measurable field depending on one coordinate
\cite[Theorem~1]{BatemanThiele2013}. The one-coordinate \(L^2\) case is already connected to
Carleson's theorem by a partial Fourier transform
\cite{Carleson1966,BatemanThiele2013}; Bateman and Thiele obtain the stated wider
\(p\)-range. The one-coordinate hypothesis remains a substantive
restriction on the direction field. Guo extended the one-coordinate setting to
fields constant along suitable Lipschitz foliations, with a uniform
transversality assumption \cite{Guo2015,Guo2017}.

A different family of results combines Lipschitz control with
lacunary directions. Guo and Thiele proved strong \(L^p\) bounds
without a frequency restriction for short principal-value transforms
associated with dyadic roundings of a positive Lipschitz slope
\cite{GuoThiele2017}; Di Plinio and Parissis extended this framework
to finite lacunary order for fields generated by separated dyadic
roundings of Lipschitz angular functions \cite{DiPlinioParissis2018}.
The rounded direction field can itself be discontinuous. These
results exploit directional structure beyond that assumed in the
present theorem.

A further reduction separates frequency reassembly from the main
single-band estimate. Di Plinio, Guo, Thiele and Zorin-Kranich
established that uniform strong \(L^2\) bounds on vertical frequency
bands imply the full short principal-value bound under a small uniform
Lipschitz bound in the vertical variable
\cite[Corollaries~1.3--1.4 of the author version]{DiPlinioGuoThieleZorinKranich2018}.
We give a commutator proof adapted to the two truncation endpoints
and to the narrow frequency bins used later. The new estimate to be
proved is the logarithmically improved restricted bound on a band;
the principle of band-to-full reassembly is already present in that
work.

\subsection{The analytic and geometric reductions}

We describe the proof before fixing its technical parameters.
Localization and rotation write a short piece of the operator in
slope form
\[
 H_{u,a}^{\epsilon}h(x,y)
 =\int_{\epsilon<|t|<a}
 h(x-t,y-tu(x,y))\,\frac{\mathrm dt}{t},
 \qquad |u|\le1,\quad \Lip(u)\le C,
\]
with both input and test supported in a short horizontal strip.
Writing \(\eta\) for the frequency dual to \(y\), a vertical
frequency band at scale \(w\) means that \(|\eta|\) is comparable
to \(w^{-1}\). For its smooth projection \(P_w\), the
central restricted estimate is
\[
 |\langle H_{u,a}^{\epsilon}P_w f,g\rangle|
 \lesssim
 \bigl(1+|\log(|F|/|E|)|\bigr)^{-3}\sqrt{|F||E|},
 \qquad |f|\le\one_F,\quad |g|\le\one_E.
\]
Here \(F,E\) are sets of positive finite measure in the horizontal
strip. Its constants are uniform in \(w\) and \(\epsilon\). The logarithmic
factor makes amplitude-level summation converge. A Lipschitz
commutator then permits Littlewood--Paley reassembly across vertical
bands. These implications, including all endpoint errors, are proved
in this section before the restricted estimate itself is used.

The difficult regime is \(|E|\gg|F|\). Write
\(L=\log\sqrt{|E|/|F|}\). Positive operator bounds discard all but
very short integration times. The vertical band is then divided into
narrow intervals centered at frequencies \(\eta_j\), with common
width comparable to \(W^{-1}\), where \(W=mw\) and \(m\) grows
subexponentially in \(L\). The components form one finite-dimensional
Hilbert space. A wave packet has horizontal length \(\ell\),
transverse width \(W\), and a direction label with precision
\(d=w/\ell\). Smooth packet forms retain the original field only
through scalar selectors of its direction.

Ordinary size and density estimates remove the small-density packets.
Each remaining packet is assigned to a top box whose direction is
popular on a measured subset of the test set. The popularity witnesses
are disjoint. Lacey--Li's theorem controls a decaying weighted count
of the tops, after every thin parallelogram has been enclosed by an
admissible orthogonal rectangle. This step is geometric: it supplies
a polynomial cost in \(L\), and does not provide the final
oscillatory saving.

\subsection{The phase gain and its iteration}

Two statements make the oscillation usable. First, on every reference
line the packet sums have polynomial moment bounds, square-function
bounds and \(3\)-variation bounds. The variation argument combines
L\'epingle's martingale inequality with the comparison between
smooth Fourier truncations and conditional expectations developed
by Jones, Seeger and Wright \cite{Lepingle1976,JSW2008}. All three
packet estimates hold for every fixed deterministic subcollection. Popularity and Lipschitz control then
produce one exceptional set outside which each separated list of
significant top directions in the required angular cluster and
precision range has at most \(D=L^C\) members, simultaneously over
angular resolutions.

Second, a finite Fourier sum with at most \(D\) separated real
frequencies has a strict improvement over the usual coefficient
exponent. For suitable band-limited probability measures \(\nu\),
we prove
\[
 \left\|\sum_{i=1}^d a_i e^{i\lambda_i t}\right\|_{L^p(\nu)}
 \le \left(\sum_{i=1}^d |a_i|^{r_*}\right)^{1/r_*},
 \qquad
 \frac1{r_*}=1-\frac{1+c/\log(2D)}p,
 \quad p\ge4.
\]
The operator constant is exactly one, for complex coefficients and
for Hilbert-valued coefficients acted on by diagonal phases.
The proof separates the case of one dominant coefficient from its
complement. In the latter case, separated quantiles give a strict
fourth-moment deficit, which survives a small improvement of the
coefficient exponent. An unspecified constant in this inequality
would accumulate at every angular level and destroy the argument.

The iteration organizes directions into nested angular intervals.
Each interval carries the positive mass of the tops whose direction
precision is sufficiently fine for that interval; these are its
eligible tops. It also carries an oscillatory amplitude assembled
from its packets. The significant-direction count enters phase
averaging through a weighted selection of separated block children,
whose analytic weights cancel the probability cost of retaining them.
We compare mass-normalized \(p\)-energies after adding the mass itself
as a baseline. On an appropriate spatial interval, a finer amplitude
can be predicted from one good point by diagonal phase motion. The
finite Fourier estimate therefore controls the average of the
predicted energy factor. The comparison errors are proportional to the
relative mass lost on the path, so their sum is controlled by a
logarithmic mass telescope.

Packets that stop between two angular endpoints create residual
terms. Square summation handles blocks with a definite mass drop.
Across nearly unary portions of a path, the energies telescope and
the remaining residual sums are estimated by \(3\)-variation of
one fixed sequence for that path. The sequence is chosen before the
spatial point or mass stop is selected. Normalized factors on nested
spatial cells can then be integrated without any cost proportional
to the angular depth.

Finally we interpolate among shifted groupings of angular levels.
At the central parameter the packet coefficients contain boundary
weights. Averaging one common angular-grid translation removes these
weights, using eligibility as well as angular proximity in the
boundary telescope. With \(p=L^{1/2}\), the strict phase gain leaves
a decrease of order \((\sqrt L\log L)^{-1}\) in the exponent of
\(\sqrt{|E|/|F|}\). This dominates all fixed polynomial losses and
supplies the logarithmic restricted estimate above.

\subsection{Organization, inputs and uniformity}

Section~\ref{intro:section} gives the analytic reduction and the
precise packet estimates needed for the band theorem.
Section~\ref{tile:section} proves the ordinary tile estimates and
constructs the popular forest.
Section~\ref{conc:section} establishes every-line concentration and
the simultaneous count of significant directions.
Section~\ref{phase:section} proves the constant-one finite Fourier
inequality.
Section~\ref{hier:section} constructs the angular hierarchy and the
predicted factors.
Section~\ref{res:section} controls residuals and integrates path
products.
Section~\ref{assembly:section} performs interpolation, angular-grid
averaging and the final strong-type assembly.

There are two specialized external inputs: Lacey--Li's short
Lipschitz-popularity maximal theorem, stated in
\Cref{red:kakeya}, and the unweighted scalar martingale
\(3\)-variation inequality of L\'epingle, used in the form
\cite[Theorem~1.1, \(p=2\), \(r=3\), weight one]{ZorinKranich2020}.
The latter enters the locally proved Fourier/martingale comparison
in \Cref{conc:annular-variation}; we credit its origin to
\cite{Lepingle1976}. Standard one-dimensional Hilbert-transform,
maximal-function and interpolation facts are recalled where used.
The commutator, tile, concentration, finite-phase and hierarchy
arguments are included in full.

All packet collections, scale ranges and angular depths are finite
until estimates independent of those truncations have been proved.
The bin Hilbert spaces are finite-dimensional and complex, and every
constant is independent of their dimension. Pairings are linear in
the first variable. A factor \(L^C\) always has a fixed exponent;
it never conceals the number of packets, scales or hierarchy levels.
